\documentclass[10pt]{article}
 
\usepackage[utf8]{inputenc}
\usepackage[T1]{fontenc}
\usepackage[english]{babel}
\usepackage{amsmath, amssymb}
\usepackage{hyperref}
%\usepackage[margin=2cm]{geometry}
\usepackage[numbers,sort&compress]{natbib}
%\usepackage{microtype}
 
%\setlength{\parskip}{4pt}
%\setlength{\bibsep}{1pt plus 0.3ex}
 
\title{\vspace{-1.5em}\textbf{Characterising functions for which neural network depth is
strictly necessary}\vspace{-0.5em}}
\author{Romain Ben \quad Thibaud Crotta \quad Zouhair Saitout \\[2pt]
\small Polytech Nice Sophia, MAM3}
\date{}
 
\begin{document}
\maketitle

Deep neural networks owe much of their empirical superiority to depth: deeper architectures
consistently outperform wider but shallower ones at equal parameter count, suggesting that depth
is a fundamental source of expressive power rather than a contingent design choice. This stands
in tension with classical universality results~\citep{cybenko89,hornik91}, which guarantee that
a single hidden layer suffices to approximate any continuous function on a compact set. The
paradox dissolves once approximation \emph{cost} is considered: for some targets, every shallow
network must be exponentially wider than a deep one to reach the same accuracy. This phenomenon
is called \emph{depth separation}, and the question we investigate is whether the class of
functions exhibiting it admits a mathematical characterisation beyond isolated examples.

The universality results of \citet{cybenko89} and \citet{hornik91}, surveyed by
\citet{pinkus99}, give no control on the required width~$N$, leaving open an exponential
dependence on complexity or input dimension~$d$. \citet{barron93} obtained a dimension-free
bound for functions whose Fourier transform satisfies $C_f := \int\|\omega\|\,|\hat
f(\omega)|\,\mathrm{d}\omega < \infty$: a sigmoidal network with $N$ neurons then achieves
squared $L^2$ error at most $C_f^2/N$. For more general functions, \citet{yarotsky17} showed
that deep ReLU networks attain near-optimal rates over Sobolev balls $W^{k,\infty}([0,1]^d)$
with depth scaling as $\log(1/\varepsilon)$, indicating that depth becomes essential for
smooth-but-rough targets.

Several families of provable separations confirm this picture. \citet{telgarsky16} constructs
the iterated triangle wave $t_k$ and proves that any shallow ReLU network approximating $t_k$
to constant accuracy requires $2^{k/6}$ neurons, while depth~$k$ suffices with $O(k)$ neurons
,a gap rooted in the fact that each layer at most doubles the number of linear pieces.
\citet{montufar14} show that the number of linear regions grows exponentially in depth but only
polynomially in width, confirming depth as the decisive factor for geometric expressivity.
\citet{eldan16} sharpen the separation to a single layer boundary: a radial function
approximable by a depth-$3$ network of polynomial size requires $\exp(\Omega(d))$ neurons at
depth~$2$, with the obstruction located in the radial Fourier spectrum. \citet{safran17} extend
this trade-off to natural targets such as $f(x)=\|x\|_2$.

The search for a general characterisation has produced partial but precise results.
\citet{daniely17} gives an algebraic sufficient condition for non-approximability by shallow
networks, while \citet{raghu17} formalise expressivity via the exponential growth of curve
lengths with depth. The most complete result is due to \citet{diakonikolas22}: on
$\mathbb{S}^{d-1}$, a function is approximable by a polynomial-size depth-$2$ network if and
only if its spherical Fourier spectrum concentrates on harmonics of degree
$O(\mathrm{polylog}\,d)$, the only known necessary and sufficient condition, though
restricted to the sphere and to $L^2$.

A fundamental obstacle is identified by \citet{vardi20}: every known separation relies on
polynomially bounded weights, and with unrestricted weights a shallow network can simulate any
deep one exactly, collapsing all known separations. Any future characterisation must therefore
specify this regime explicitly. Two broader gaps remain: extending the spectral criterion of
\citet{diakonikolas22} to general compact domains, and formalising the connection with Boolean
circuit complexity, where depth hierarchies such as $\mathrm{AC}^0\subsetneq\mathrm{NC}^1$ are
well understood and could, if transferred, ground architectural design in rigorous theory.

\bibliographystyle{abbrvnat}
\bibliography{biblio}

\end{document}